% Options for packages loaded elsewhere
\PassOptionsToPackage{unicode}{hyperref}
\PassOptionsToPackage{hyphens}{url}
\PassOptionsToPackage{dvipsnames,svgnames,x11names}{xcolor}
\documentclass[
  french,
  10pt,
  a4paper,
]{article}
\usepackage{xcolor}
\usepackage[margin=2.2cm]{geometry}
\usepackage{amsmath,amssymb}
\setcounter{secnumdepth}{-\maxdimen} % remove section numbering
\usepackage{iftex}
\ifPDFTeX
  \usepackage[T1]{fontenc}
  \usepackage[utf8]{inputenc}
  \usepackage{textcomp} % provide euro and other symbols
\else % if luatex or xetex
  \usepackage{unicode-math} % this also loads fontspec
  \defaultfontfeatures{Scale=MatchLowercase}
  \defaultfontfeatures[\rmfamily]{Ligatures=TeX,Scale=1}
\fi
\usepackage{lmodern}
\ifPDFTeX\else
  % xetex/luatex font selection
\fi
% Use upquote if available, for straight quotes in verbatim environments
\IfFileExists{upquote.sty}{\usepackage{upquote}}{}
\IfFileExists{microtype.sty}{% use microtype if available
  \usepackage[]{microtype}
  \UseMicrotypeSet[protrusion]{basicmath} % disable protrusion for tt fonts
}{}
\makeatletter
\@ifundefined{KOMAClassName}{% if non-KOMA class
  \IfFileExists{parskip.sty}{%
    \usepackage{parskip}
  }{% else
    \setlength{\parindent}{0pt}
    \setlength{\parskip}{6pt plus 2pt minus 1pt}}
}{% if KOMA class
  \KOMAoptions{parskip=half}}
\makeatother
\usepackage{longtable,booktabs,array}
\usepackage{calc} % for calculating minipage widths
% Correct order of tables after \paragraph or \subparagraph
\usepackage{etoolbox}
\makeatletter
\patchcmd\longtable{\par}{\if@noskipsec\mbox{}\fi\par}{}{}
\makeatother
% Allow footnotes in longtable head/foot
\IfFileExists{footnotehyper.sty}{\usepackage{footnotehyper}}{\usepackage{footnote}}
\makesavenoteenv{longtable}
\ifLuaTeX
\usepackage[bidi=basic,shorthands=off,]{babel}
\else
\usepackage[bidi=default,shorthands=off,]{babel}
\fi
\ifLuaTeX
  \usepackage{selnolig} % disable illegal ligatures
\fi
\setlength{\emergencystretch}{3em} % prevent overfull lines
\providecommand{\tightlist}{%
  \setlength{\itemsep}{0pt}\setlength{\parskip}{0pt}}
% Fichier LaTeX généré par docs/build/build_pdf.sh à partir de 03_conception_MVP.md.
% Compilation (depuis le dossier de ce fichier), deux fois pour la table des matières :
%   pdflatex 04_Guide3_Conception_MVP.tex
% Pour insérer une capture : déposez le fichier PNG au nom indiqué dans le cadre « CAPTURE À INSÉRER »
% (dossier ../images/...), puis recompilez. Vous pouvez aussi remplacer la commande \CredixCapture{...}{...}
% par votre propre \includegraphics.
\newcommand{\CredixShortTitle}{Guide 3 : MVP}
% Préambule commun aux PDF du dossier de remise CREDIX (compilé avec pdfLaTeX via pandoc)
\usepackage[scaled=.92]{helvet}
\renewcommand{\familydefault}{\sfdefault}
\usepackage{pifont}
\usepackage{amssymb}
% Caractères Unicode utilisés dans les documents (pdfLaTeX ne les connaît pas d'office)
\DeclareUnicodeCharacter{03C1}{\ensuremath{\rho}}
\DeclareUnicodeCharacter{2265}{\ensuremath{\geq}}
\DeclareUnicodeCharacter{2264}{\ensuremath{\leq}}
\DeclareUnicodeCharacter{2248}{\ensuremath{\approx}}
\DeclareUnicodeCharacter{2192}{\ensuremath{\rightarrow}}
\DeclareUnicodeCharacter{2190}{\ensuremath{\leftarrow}}
\DeclareUnicodeCharacter{2212}{\ensuremath{-}}
\DeclareUnicodeCharacter{2713}{\ding{51}}
\DeclareUnicodeCharacter{2714}{\ding{52}}
\DeclareUnicodeCharacter{2717}{\ding{55}}
\DeclareUnicodeCharacter{25B2}{\ensuremath{\blacktriangle}}
\DeclareUnicodeCharacter{25CB}{\ensuremath{\circ}}
\DeclareUnicodeCharacter{2500}{-}
\DeclareUnicodeCharacter{2502}{|}
\DeclareUnicodeCharacter{250C}{+}
\DeclareUnicodeCharacter{2510}{+}
\DeclareUnicodeCharacter{2514}{+}
\DeclareUnicodeCharacter{2518}{+}
\DeclareUnicodeCharacter{251C}{+}
\DeclareUnicodeCharacter{2524}{+}
\DeclareUnicodeCharacter{252C}{+}
\DeclareUnicodeCharacter{2534}{+}
\DeclareUnicodeCharacter{253C}{+}
\DeclareUnicodeCharacter{0192}{\textit{f}}
\DeclareUnicodeCharacter{2011}{-}
\DeclareUnicodeCharacter{202F}{\,}
\DeclareUnicodeCharacter{2009}{\,}
\DeclareUnicodeCharacter{00D7}{\ensuremath{\times}}
\DeclareUnicodeCharacter{25A0}{\ensuremath{\blacksquare}}
\DeclareUnicodeCharacter{2022}{\ensuremath{\bullet}}
\DeclareUnicodeCharacter{2026}{\ldots}
\DeclareUnicodeCharacter{2460}{\ding{172}}
\DeclareUnicodeCharacter{2461}{\ding{173}}
\DeclareUnicodeCharacter{2462}{\ding{174}}
\DeclareUnicodeCharacter{2463}{\ding{175}}
\DeclareUnicodeCharacter{2464}{\ding{176}}
\DeclareUnicodeCharacter{2465}{\ding{177}}
\DeclareUnicodeCharacter{2466}{\ding{178}}
\DeclareUnicodeCharacter{2467}{\ding{179}}
\usepackage{fvextra}
\DefineVerbatimEnvironment{verbatim}{Verbatim}{breaklines=true,breakanywhere=true,fontsize=\footnotesize,frame=single,framerule=0.3pt,rulecolor=\color[gray]{0.7},framesep=2mm,xleftmargin=1mm,xrightmargin=1mm}
\fvset{breaklines=true,breakanywhere=true}
\usepackage{fancyhdr}
\usepackage{tcolorbox}
\usepackage{colortbl}
\usepackage{enumitem}
\definecolor{credixblue}{HTML}{1F3A5F}
\definecolor{credixgray}{HTML}{5F6B7A}
\definecolor{boxbg}{HTML}{F3F5F8}
\setlength{\emergencystretch}{4em}
\renewcommand{\arraystretch}{1.3}
\setlength{\parindent}{0pt}
\setlength{\parskip}{0.55em}
\setlist{itemsep=0.15em,topsep=0.3em}
\pagestyle{fancy}
\fancyhf{}
\fancyhead[L]{\small\color{credixgray}\textsc{\CredixShortTitle}}
\fancyhead[R]{\small\color{credixgray}CREDIX --- Dossier de remise}
\fancyfoot[C]{\small\color{credixgray}\thepage}
\renewcommand{\headrulewidth}{0.3pt}
\renewcommand{\footrulewidth}{0pt}
\fancypagestyle{plain}{\fancyhf{}\fancyfoot[C]{\small\color{credixgray}\thepage}\renewcommand{\headrulewidth}{0pt}}
\usepackage{titlesec}
\titleformat{\section}{\Large\bfseries\color{credixblue}}{}{0em}{}[\vspace{-0.2em}\color{credixblue}\titlerule]
\titleformat{\subsection}{\large\bfseries\color{credixblue}}{}{0em}{}
\titleformat{\subsubsection}{\normalsize\bfseries\color{credixblue}}{}{0em}{}
% Cadre de repère pour une capture manquante
\newcommand{\CaptureManquante}[2]{%
  \begin{tcolorbox}[colback=boxbg,colframe=credixgray,boxrule=0.5pt,arc=1mm,left=3mm,right=3mm,top=3mm,bottom=3mm,halign=center]
  \textbf{CAPTURE À INSÉRER}\\[0.4em]#1\\[0.4em]{\footnotesize\ttfamily #2}
  \end{tcolorbox}}

% Images : insérées si le fichier existe au moment de la compilation, sinon cadre « CAPTURE À INSÉRER ».
% Pour ajouter une capture : déposer le PNG au nom indiqué dans le cadre, puis recompiler.
\usepackage{graphicx}
\newcommand{\CredixCapture}[2]{%
  \IfFileExists{#2}{%
    \begin{center}\includegraphics[width=\linewidth,height=0.75\textheight,keepaspectratio]{#2}\end{center}}%
  {\CaptureManquante{#1}{\detokenize{#2}}}}
\newcommand{\CredixFigure}[1]{%
  \begin{center}\includegraphics[width=0.95\linewidth,height=0.8\textheight,keepaspectratio]{#1}\end{center}}
\usepackage{bookmark}
\IfFileExists{xurl.sty}{\usepackage{xurl}}{} % add URL line breaks if available
\urlstyle{same}
\hypersetup{
  pdftitle={Guide 3 : conception du MVP},
  pdfauthor={SINGHE PENKA Hendrix Donavan},
  pdflang={fr},
  colorlinks=true,
  linkcolor={credixblue},
  filecolor={Maroon},
  citecolor={Blue},
  urlcolor={credixblue},
  pdfcreator={LaTeX via pandoc}}

\author{}
\date{}

\begin{document}

\begin{titlepage}
\thispagestyle{empty}
\centering
\vspace*{3.2cm}
{\color{credixgray}\large\textsc{Projet de fin d'études --- ENSPY, Université de Yaoundé I}\par}
\vspace{2.2cm}
{\Huge\bfseries\color{credixblue} Guide 3 : conception du MVP\par}
\vspace{0.8cm}
{\Large Un scoring de crédit réduit, à construire avec Claude Code\par}
\vspace{2cm}
{\color{credixblue}\rule{0.5\linewidth}{0.8pt}\par}
\vspace{1.2cm}
{\large SINGHE PENKA Hendrix Donavan\par}
{\color{credixgray}Dépôt GitHub : HendrixPenka2/credix\par}
\vfill
{\color{credixgray}\small Version du 19/09/2026\par}
\end{titlepage}


\renewcommand*\contentsname{Table des matières}
{
\setcounter{tocdepth}{2}
\tableofcontents
}
\section{CREDIX-MVP : document de
conception}\label{credix-mvp--document-de-conception}

\textbf{Un mini-système de scoring de crédit, conçu pour être construit
avec Claude Code en suivant la méthode du guide 2.}

\begin{center}\rule{0.5\linewidth}{0.5pt}\end{center}

\subsection{1. Introduction}\label{1-introduction}

\subsubsection{1.1 Contexte}\label{11-contexte}

\textbf{CREDIX} est un système complet de scoring de crédit (voir le
README du dépôt). Il compte plusieurs dizaines de routes, deux modèles
complémentaires, une base documentaire, un service de rapports PDF et
trois espaces utilisateurs. Il est \textbf{trop grand} pour servir
d\textquotesingle exercice à une démonstration.

\textbf{CREDIX-MVP} en est une version \textbf{réduite} : le plus petit
système qui garde l\textquotesingle idée essentielle. \emph{Un agent
saisit une demande de crédit ; le système calcule une probabilité de
défaut, la convertit en score, propose une décision et explique les
trois facteurs qui pèsent le plus.}

Ce document est écrit pour être donné à \textbf{Claude Code} en mode
Plan (guide 1, section 8) : il contient tout ce qu\textquotesingle il
faut pour produire un plan de réalisation, puis un système qui tourne.

\subsubsection{1.2 Objectifs}\label{12-objectifs}

\begin{longtable}[]{@{}
  >{\raggedright\arraybackslash}p{(\linewidth - 2\tabcolsep) * \real{0.1105}}
  >{\raggedright\arraybackslash}p{(\linewidth - 2\tabcolsep) * \real{0.8595}}@{}}
\toprule\noalign{}
\begin{minipage}[b]{\linewidth}\raggedright
N°
\end{minipage} & \begin{minipage}[b]{\linewidth}\raggedright
Objectif
\end{minipage} \\
\midrule\noalign{}
\endhead
\bottomrule\noalign{}
\endlastfoot
\textbf{O1} & Disposer d\textquotesingle un système de scoring
\textbf{de bout en bout} (données, modèle, API, interface, tests)
réalisable en quelques modules \\
\textbf{O2} & \textbf{Valider la méthode} du guide 2 : à partir de ce
seul document, Claude Code doit produire un plan cohérent, puis un
premier module qui passe ses tests \\
\textbf{O3} & Garder les \textbf{principes de CREDIX} : score de 300 à
850, décision à trois issues, explication des facteurs, traçabilité,
seuils réglables \\
\end{longtable}

\subsubsection{1.3 Périmètre}\label{13-puxe9rimuxe8tre}

\begin{longtable}[]{@{}
  >{\raggedright\arraybackslash}p{(\linewidth - 2\tabcolsep) * \real{0.4942}}
  >{\raggedright\arraybackslash}p{(\linewidth - 2\tabcolsep) * \real{0.4758}}@{}}
\toprule\noalign{}
\begin{minipage}[b]{\linewidth}\raggedright
Inclus dans le MVP
\end{minipage} & \begin{minipage}[b]{\linewidth}\raggedright
Exclu (voir la section 13)
\end{minipage} \\
\midrule\noalign{}
\endhead
\bottomrule\noalign{}
\endlastfoot
Deux rôles : \textbf{Agent} et \textbf{Administrateur} & Rôle
\textbf{Superviseur} et file de revue manuelle \\
Un modèle \textbf{statistique simple} (régression logistique) & Deux
flux, autoencodeur, calibration isotonique, WoE \\
Six variables explicatives & Vingt-sept variables et indice de
couverture ρc \\
Base \textbf{SQLite}, aucun Docker & MongoDB, conteneurs, service PDF \\
Une page web unique & Trois espaces, thème sombre, graphiques \\
Données \textbf{synthétiques} générées par un script & Jeu de données
réel \\
\end{longtable}

\subsubsection{1.4 Comment utiliser ce document avec Claude
Code}\label{14-comment-utiliser-ce-document-avec-claude-code}

\begin{enumerate}
\def\labelenumi{\arabic{enumi}.}
\tightlist
\item
  Créez un dossier vide et ouvrez-le dans \textbf{VS Code}.
\item
  Copiez-y ce fichier sous
  \texttt{docs/\allowbreak{}conception\_\allowbreak{}mvp.\allowbreak{}md}.
\item
  Ouvrez Claude Code et passez en \textbf{mode Plan}.
\item
  Donnez le prompt de l\textquotesingle{}\textbf{annexe B}.
\end{enumerate}

Le déroulé complet de la démonstration est dans la fiche de
démonstration.

\begin{center}\rule{0.5\linewidth}{0.5pt}\end{center}

\subsection{2. Exigences}\label{2-exigences}

\subsubsection{2.1 Exigences
fonctionnelles}\label{21-exigences-fonctionnelles}

\begin{longtable}[]{@{}
  >{\raggedright\arraybackslash}p{(\linewidth - 4\tabcolsep) * \real{0.1125}}
  >{\raggedright\arraybackslash}p{(\linewidth - 4\tabcolsep) * \real{0.4287}}
  >{\raggedright\arraybackslash}p{(\linewidth - 4\tabcolsep) * \real{0.4287}}@{}}
\toprule\noalign{}
\begin{minipage}[b]{\linewidth}\raggedright
N°
\end{minipage} & \begin{minipage}[b]{\linewidth}\raggedright
Exigence
\end{minipage} & \begin{minipage}[b]{\linewidth}\raggedright
Critère d\textquotesingle acceptation
\end{minipage} \\
\midrule\noalign{}
\endhead
\bottomrule\noalign{}
\endlastfoot
\textbf{EF-1} & Un utilisateur se \textbf{connecte} avec un identifiant
et un mot de passe et reçoit un \textbf{jeton} valable 8 heures & Bons
identifiants : jeton reçu ; mauvais identifiants : refus (401) \\
\textbf{EF-2} & L\textquotesingle accès aux routes dépend du
\textbf{rôle} (Agent, Administrateur) & Un Agent qui appelle une route
d\textquotesingle administration reçoit 403 \\
\textbf{EF-3} & Un agent \textbf{crée un client} (nom, prénom, âge,
revenu mensuel, ancienneté d\textquotesingle emploi, nombre
d\textquotesingle incidents de paiement) & Le client créé est
retrouvable ; une valeur invalide (âge négatif) est refusée (422) \\
\textbf{EF-4} & Un agent \textbf{recherche} un client par nom ou prénom
et \textbf{consulte} sa fiche & La recherche « ndi » retrouve « Ndiaye
» \\
\textbf{EF-5} & Un agent \textbf{demande un scoring} pour un client,
avec un montant et une durée de crédit & Le résultat contient une
probabilité de défaut, un \textbf{score entre 300 et 850} et une
\textbf{décision} \\
\textbf{EF-6} & Le résultat contient \textbf{trois facteurs} qui
expliquent le score (libellé, contribution, sens) & Trois facteurs,
triés par importance décroissante \\
\textbf{EF-7} & Un agent consulte l\textquotesingle{}\textbf{historique
des scorings} d\textquotesingle un client & Les scorings sont listés du
plus récent au plus ancien \\
\textbf{EF-8} & Un administrateur \textbf{crée un compte} (agent ou
administrateur) & Le nouveau compte peut se connecter \\
\textbf{EF-9} & Un administrateur \textbf{règle les deux seuils} de
décision & Un seuil incohérent (accord supérieur au refus) est refusé
(422) \\
\textbf{EF-10} & Un script \textbf{génère les données} et
\textbf{entraîne le modèle}, puis affiche ses performances & Le script
écrit le modèle et affiche une AUC d\textquotesingle au moins 0,70 \\
\end{longtable}

\subsubsection{2.2 Exigences non
fonctionnelles}\label{22-exigences-non-fonctionnelles}

\begin{longtable}[]{@{}
  >{\raggedright\arraybackslash}p{(\linewidth - 4\tabcolsep) * \real{0.1125}}
  >{\raggedright\arraybackslash}p{(\linewidth - 4\tabcolsep) * \real{0.4547}}
  >{\raggedright\arraybackslash}p{(\linewidth - 4\tabcolsep) * \real{0.4028}}@{}}
\toprule\noalign{}
\begin{minipage}[b]{\linewidth}\raggedright
N°
\end{minipage} & \begin{minipage}[b]{\linewidth}\raggedright
Exigence
\end{minipage} & \begin{minipage}[b]{\linewidth}\raggedright
Critère d\textquotesingle acceptation
\end{minipage} \\
\midrule\noalign{}
\endhead
\bottomrule\noalign{}
\endlastfoot
\textbf{ENF-1} & \textbf{Traçabilité} : chaque scoring garde la version
du modèle, les seuils appliqués, l\textquotesingle utilisateur et la
date & Ces quatre informations figurent dans
l\textquotesingle historique \\
\textbf{ENF-2} & \textbf{Sécurité} : mots de passe hachés, jamais en
clair ; secret du jeton dans une variable
d\textquotesingle environnement & Aucun mot de passe lisible dans la
base ni dans le code \\
\textbf{ENF-3} & \textbf{Reproductibilité} : deux exécutions du script
avec la même graine donnent le même modèle & Mêmes coefficients à la
troisième décimale \\
\textbf{ENF-4} & \textbf{Lancement simple} : installation et démarrage
en six commandes, sans Docker & Le README est suivi à la lettre sur une
machine neuve \\
\textbf{ENF-5} & \textbf{Réactivité} : le modèle est chargé une seule
fois au démarrage & Un scoring répond en moins de 500 ms en local \\
\textbf{ENF-6} & \textbf{Testabilité} : une commande unique lance tous
les tests & \texttt{pytest} passe, avec au moins un test par exigence
fonctionnelle \\
\end{longtable}

\begin{center}\rule{0.5\linewidth}{0.5pt}\end{center}

\subsection{3. Acteurs et cas
d\textquotesingle utilisation}\label{3-acteurs-et-cas-dutilisation}

Le système a \textbf{deux acteurs}.
L\textquotesingle{}\textbf{Administrateur} peut faire tout ce que fait
l\textquotesingle Agent, et en plus gérer les comptes et les seuils : il
\textbf{hérite} du rôle Agent.

La figure 1 présente les cas d\textquotesingle utilisation de chaque
acteur.

\CredixFigure{../images/mermaid/03_conception_MVP_1.png}

\emph{Figure 1. Cas d\textquotesingle utilisation de CREDIX-MVP. Les
ovales sont les cas d\textquotesingle utilisation, les cercles les
acteurs. La flèche en pointillés indique que
l\textquotesingle Administrateur hérite de tout ce que fait
l\textquotesingle Agent.}

\subsubsection{3.1 Spécification du cas « Demander un scoring
»}\label{31-spuxe9cification-du-cas--demander-un-scoring-}

\begin{longtable}[]{@{}
  >{\raggedright\arraybackslash}p{(\linewidth - 2\tabcolsep) * \real{0.2250}}
  >{\raggedright\arraybackslash}p{(\linewidth - 2\tabcolsep) * \real{0.7450}}@{}}
\toprule\noalign{}
\begin{minipage}[b]{\linewidth}\raggedright
Élément
\end{minipage} & \begin{minipage}[b]{\linewidth}\raggedright
Contenu
\end{minipage} \\
\midrule\noalign{}
\endhead
\bottomrule\noalign{}
\endlastfoot
\textbf{Acteur} & Agent (ou Administrateur) \\
\textbf{Objectif} & Obtenir une décision expliquée pour une demande de
crédit \\
\textbf{Préconditions} & L\textquotesingle agent est connecté ; le
client existe \\
\textbf{Scénario nominal} & 1. L\textquotesingle agent ouvre la fiche du
client. 2. Il saisit le \textbf{montant} et la \textbf{durée} du crédit.
3. Il valide. 4. Le système calcule la probabilité de défaut, le score
et la décision, et sélectionne les trois facteurs principaux. 5. Le
système enregistre le scoring. 6. Il affiche le résultat. \\
\textbf{Alternatives} & \textbf{3a.} Montant ou durée invalides : le
système refuse (422) et l\textquotesingle indique. \textbf{4a.} Le
modèle n\textquotesingle est pas disponible : erreur 503 avec un message
clair. \\
\textbf{Postcondition} & Un scoring est enregistré avec sa version de
modèle, ses seuils, son auteur et sa date \\
\end{longtable}

\subsubsection{3.2 Spécification du cas « Créer un compte
»}\label{32-spuxe9cification-du-cas--cruxe9er-un-compte-}

\begin{longtable}[]{@{}
  >{\raggedright\arraybackslash}p{(\linewidth - 2\tabcolsep) * \real{0.2250}}
  >{\raggedright\arraybackslash}p{(\linewidth - 2\tabcolsep) * \real{0.7450}}@{}}
\toprule\noalign{}
\begin{minipage}[b]{\linewidth}\raggedright
Élément
\end{minipage} & \begin{minipage}[b]{\linewidth}\raggedright
Contenu
\end{minipage} \\
\midrule\noalign{}
\endhead
\bottomrule\noalign{}
\endlastfoot
\textbf{Acteur} & Administrateur \\
\textbf{Préconditions} & L\textquotesingle administrateur est
connecté \\
\textbf{Scénario nominal} & 1. L\textquotesingle administrateur saisit
un identifiant, un mot de passe (8 caractères minimum) et un rôle. 2. Le
système vérifie que l\textquotesingle identifiant est libre. 3. Il
enregistre le compte, mot de passe \textbf{haché}. \\
\textbf{Alternatives} & \textbf{2a.} Identifiant déjà pris : refus
(409). \textbf{1a.} Mot de passe trop court : refus (422). \\
\end{longtable}

\subsubsection{3.3 Spécification du cas « Régler les seuils de décision
»}\label{33-spuxe9cification-du-cas--ruxe9gler-les-seuils-de-duxe9cision-}

\begin{longtable}[]{@{}
  >{\raggedright\arraybackslash}p{(\linewidth - 2\tabcolsep) * \real{0.1895}}
  >{\raggedright\arraybackslash}p{(\linewidth - 2\tabcolsep) * \real{0.7805}}@{}}
\toprule\noalign{}
\begin{minipage}[b]{\linewidth}\raggedright
Élément
\end{minipage} & \begin{minipage}[b]{\linewidth}\raggedright
Contenu
\end{minipage} \\
\midrule\noalign{}
\endhead
\bottomrule\noalign{}
\endlastfoot
\textbf{Acteur} & Administrateur \\
\textbf{Scénario nominal} & 1. L\textquotesingle administrateur saisit
un \textbf{seuil d\textquotesingle accord} et un \textbf{seuil de refus}
(probabilités de défaut). 2. Le système vérifie que 0 \textless{} accord
\textless{} refus \textless{} 1. 3. Il enregistre une \textbf{nouvelle
version} des seuils (l\textquotesingle historique est conservé). \\
\textbf{Règle métier} & Probabilité de défaut \textbf{inférieure au
seuil d\textquotesingle accord} : décision ACCORDÉ. \textbf{Supérieure
au seuil de refus} : REFUSÉ. \textbf{Entre les deux} : REVUE (un humain
doit trancher). Valeurs initiales : 0,10 et 0,30 (les mêmes que
CREDIX). \\
\end{longtable}

\begin{center}\rule{0.5\linewidth}{0.5pt}\end{center}

\subsection{4. Modèle du domaine}\label{4-moduxe8le-du-domaine}

Le domaine compte cinq classes principales, trois types énumérés, et les
associations qui les relient. La figure 2 présente le diagramme de
classes.

\CredixFigure{../images/mermaid/03_conception_MVP_2.png}

\emph{Figure 2. Diagramme de classes du domaine. Les losanges pleins
indiquent une composition : un Facteur n\textquotesingle existe que dans
le Scoring qu\textquotesingle il explique. Les types énumérés portent le
stéréotype « enumeration ». Les attributs sont privés (signe moins), les
opérations publiques (signe plus).}

\textbf{Règles de lecture et de conception :}

\begin{itemize}
\tightlist
\item
  Aucun attribut n\textquotesingle a pour type une autre classe du
  diagramme : les liens entre classes sont des \textbf{associations},
  avec leurs multiplicités aux deux extrémités.
\item
  Un \textbf{Scoring} garde les seuils \textbf{appliqués}
  (\texttt{seuilAccordApplique}, \texttt{seuilRefusApplique}) : ils ne
  changent pas si les seuils sont modifiés plus tard (traçabilité,
  ENF-1).
\item
  Un \textbf{Scoring} est toujours produit par \textbf{un} modèle,
  \textbf{un} utilisateur, pour \textbf{un} client.
\end{itemize}

\begin{center}\rule{0.5\linewidth}{0.5pt}\end{center}

\subsection{5. Scénario principal : demander un
scoring}\label{5-scuxe9nario-principal--demander-un-scoring}

La figure 3 montre les échanges entre l\textquotesingle agent et les
composants du système lors d\textquotesingle un scoring.

\CredixFigure{../images/mermaid/03_conception_MVP_3.png}

\emph{Figure 3. Diagramme de séquence du cas « Demander un scoring ».
Les flèches pleines sont des appels, les flèches pointillées des
retours.}

\begin{center}\rule{0.5\linewidth}{0.5pt}\end{center}

\subsection{6. Architecture et
technologies}\label{6-architecture-et-technologies}

\subsubsection{6.1 Vue
d\textquotesingle ensemble}\label{61-vue-densemble}

La figure 4 présente les composants et leurs liens.

\CredixFigure{../images/mermaid/03_conception_MVP_4.png}

\emph{Figure 4. Architecture de CREDIX-MVP. Le navigateur ne parle
qu\textquotesingle aux routes ; les services portent la logique ; le
modèle est chargé une seule fois en mémoire.}

\subsubsection{6.2 Choix technologiques}\label{62-choix-technologiques}

Chaque choix répond à une exigence.

\begin{longtable}[]{@{}
  >{\raggedright\arraybackslash}p{(\linewidth - 4\tabcolsep) * \real{0.2484}}
  >{\raggedright\arraybackslash}p{(\linewidth - 4\tabcolsep) * \real{0.4140}}
  >{\raggedright\arraybackslash}p{(\linewidth - 4\tabcolsep) * \real{0.3076}}@{}}
\toprule\noalign{}
\begin{minipage}[b]{\linewidth}\raggedright
Choix
\end{minipage} & \begin{minipage}[b]{\linewidth}\raggedright
Sert l\textquotesingle exigence
\end{minipage} & \begin{minipage}[b]{\linewidth}\raggedright
Alternative écartée
\end{minipage} \\
\midrule\noalign{}
\endhead
\bottomrule\noalign{}
\endlastfoot
\textbf{Python 3.12} & EF-10, ENF-3 : un seul langage pour
l\textquotesingle entraînement et l\textquotesingle exploitation, pas
d\textquotesingle écart entre les deux & Java : deux langages à faire
dialoguer \\
\textbf{FastAPI} & EF-2, EF-3, ENF-5 : validation automatique des
données, droits centralisés, rapidité & Flask : validation à écrire à la
main \\
\textbf{SQLite} (via SQLAlchemy) & ENF-4 : aucune installation, un
fichier & PostgreSQL, MongoDB : trop lourds pour ce MVP \\
\textbf{scikit-learn} (régression logistique) & EF-6, ENF-3 : modèle
\textbf{explicable} (contribution = coefficient × valeur standardisée),
reproductible & LightGBM : explications plus complexes (voir CREDIX) \\
\textbf{PyJWT} et \texttt{hashlib} (PBKDF2) & EF-1, ENF-2 : jeton signé,
mots de passe hachés, sans dépendance lourde & bcrypt : dépendance
native inutile ici \\
\textbf{Page HTML et JavaScript} servie par FastAPI & ENF-4 : aucun
outil de compilation & React, Next.js : chaîne de compilation en plus \\
\textbf{pytest} et \textbf{httpx} & ENF-6 & --- \\
\end{longtable}

\subsubsection{6.3 Organisation des
dossiers}\label{63-organisation-des-dossiers}

\begin{Verbatim}[breaklines=true,breakanywhere=true,fontsize=\footnotesize,frame=leftline,framerule=1.6pt,rulecolor=\color{credixblue},framesep=3mm,xleftmargin=2mm,xrightmargin=1mm]
credix-mvp/
|-- README.md
|-- requirements.txt
|-- .env.example              (SECRET_JWT)
|-- app/
|   |-- main.py               création de l'application, chargement du modèle
|   |-- config.py
|   |-- database.py           connexion SQLite, création des tables
|   |-- models.py             tables (Utilisateur, Client, Scoring, ...)
|   |-- schemas.py            formats des requêtes et des réponses
|   |-- security.py           hachage, jeton, contrôle des rôles
|   |-- routers/              auth.py, clients.py, scorings.py, admin.py
|   `-- services/             scoring.py, comptes.py, seuils.py
|-- scripts/
|   |-- generer_donnees.py    données synthétiques (annexe A)
|   |-- entrainer_modele.py   entraînement et sauvegarde du modèle
|   `-- creer_admin.py        premier compte administrateur
|-- modele/                   modele.joblib et meta.json (générés)
|-- frontend/                 index.html, app.js, style.css
`-- tests/                    un fichier de tests par module
\end{Verbatim}

\begin{center}\rule{0.5\linewidth}{0.5pt}\end{center}

\subsection{7. Modèle de données}\label{7-moduxe8le-de-donnuxe9es}

Base SQLite, une table par classe (les facteurs sont stockés avec leur
scoring).

\begin{longtable}[]{@{}
  >{\raggedright\arraybackslash}p{(\linewidth - 4\tabcolsep) * \real{0.2250}}
  >{\raggedright\arraybackslash}p{(\linewidth - 4\tabcolsep) * \real{0.3725}}
  >{\raggedright\arraybackslash}p{(\linewidth - 4\tabcolsep) * \real{0.3725}}@{}}
\toprule\noalign{}
\begin{minipage}[b]{\linewidth}\raggedright
Table
\end{minipage} & \begin{minipage}[b]{\linewidth}\raggedright
Colonnes
\end{minipage} & \begin{minipage}[b]{\linewidth}\raggedright
Contraintes
\end{minipage} \\
\midrule\noalign{}
\endhead
\bottomrule\noalign{}
\endlastfoot
\textbf{utilisateur} & \texttt{id} (clé), \texttt{identifiant},
\texttt{mot\_\allowbreak{}de\_\allowbreak{}passe\_\allowbreak{}hache},
\texttt{role}, \texttt{actif} & \texttt{identifiant} unique ;
\texttt{role} parmi AGENT, ADMIN \\
\textbf{client} & \texttt{id} (clé), \texttt{nom}, \texttt{prenom},
\texttt{age}, \texttt{revenu\_mensuel},
\texttt{anciennete\_\allowbreak{}emploi\_\allowbreak{}mois},
\texttt{nb\_\allowbreak{}incidents\_\allowbreak{}paiement},
\texttt{date\_creation}, \texttt{cree\_par} (vers utilisateur) &
\texttt{age} entre 18 et 80 ; \texttt{revenu\_mensuel} \textgreater{} 0
; ancienneté et incidents ≥ 0 \\
\textbf{modele\_scoring} & \texttt{id} (clé), \texttt{version},
\texttt{auc}, \texttt{date\_\allowbreak{}entrainement} &
\texttt{version} unique \\
\textbf{seuils} & \texttt{id} (clé), \texttt{seuil\_accord},
\texttt{seuil\_refus}, \texttt{date\_effet}, \texttt{defini\_par} (vers
utilisateur) & 0 \textless{} accord \textless{} refus \textless{} 1 \\
\textbf{scoring} & \texttt{id} (clé), \texttt{client\_id},
\texttt{utilisateur\_id}, \texttt{modele\_id}, \texttt{montant\_credit},
\texttt{duree\_mois}, \texttt{probabilite\_\allowbreak{}defaut},
\texttt{score}, \texttt{decision},
\texttt{seuil\_\allowbreak{}accord\_\allowbreak{}applique},
\texttt{seuil\_\allowbreak{}refus\_\allowbreak{}applique},
\texttt{date\_heure} & \texttt{montant\_credit} \textgreater{} 0 ;
\texttt{duree\_mois} entre 6 et 60 \\
\textbf{facteur} & \texttt{id} (clé), \texttt{scoring\_id},
\texttt{libelle}, \texttt{contribution}, \texttt{sens} & trois lignes
par scoring \\
\end{longtable}

\begin{center}\rule{0.5\linewidth}{0.5pt}\end{center}

\subsection{8. Le modèle de scoring}\label{8-le-moduxe8le-de-scoring}

\subsubsection{8.1 Variables du modèle}\label{81-variables-du-moduxe8le}

Six variables explicatives, dont deux \textbf{calculées} à partir des
données saisies :

\begin{longtable}[]{@{}
  >{\raggedright\arraybackslash}p{(\linewidth - 4\tabcolsep) * \real{0.1875}}
  >{\raggedright\arraybackslash}p{(\linewidth - 4\tabcolsep) * \real{0.5217}}
  >{\raggedright\arraybackslash}p{(\linewidth - 4\tabcolsep) * \real{0.2608}}@{}}
\toprule\noalign{}
\begin{minipage}[b]{\linewidth}\raggedright
Variable
\end{minipage} & \begin{minipage}[b]{\linewidth}\raggedright
Origine
\end{minipage} & \begin{minipage}[b]{\linewidth}\raggedright
Libellé affiché
\end{minipage} \\
\midrule\noalign{}
\endhead
\bottomrule\noalign{}
\endlastfoot
\texttt{age} & Fiche du client & « Âge » \\
\texttt{log\_revenu} & Logarithme népérien du revenu mensuel & « Revenu
mensuel » \\
\texttt{anciennete\_\allowbreak{}emploi\_\allowbreak{}mois} & Fiche du
client & « Ancienneté dans l\textquotesingle emploi » \\
\texttt{taux\_effort} & (montant du crédit ÷ durée) ÷ revenu mensuel :
\textbf{part du revenu mensuel consacrée au remboursement} & « Taux
d\textquotesingle effort » \\
\texttt{duree\_mois} & Saisie de la demande & « Durée du crédit » \\
\texttt{nb\_\allowbreak{}incidents\_\allowbreak{}paiement} & Fiche du
client & « Incidents de paiement passés » \\
\end{longtable}

\subsubsection{8.2 Entraînement}\label{82-entrauxeenement}

\begin{enumerate}
\def\labelenumi{\arabic{enumi}.}
\tightlist
\item
  Le script \texttt{generer\_\allowbreak{}donnees.\allowbreak{}py}
  produit 3 000 clients synthétiques (annexe A), avec une variable cible
  \texttt{defaut} (0 ou 1).
\item
  \texttt{entrainer\_\allowbreak{}modele.\allowbreak{}py} sépare les
  données (80 \% entraînement, 20 \% test, \textbf{graine 42}, découpage
  stratifié).
\item
  Les six variables sont \textbf{standardisées} (moyenne 0, écart-type
  1), en apprenant les moyennes et écarts-types \textbf{sur
  l\textquotesingle entraînement seulement}.
\item
  Une \textbf{régression logistique} est entraînée.
\item
  Le script affiche l\textquotesingle{}\textbf{AUC sur le jeu de test}
  et sauvegarde le modèle, les moyennes, les écarts-types, les
  coefficients et l\textquotesingle AUC dans \texttt{modele/}.
\end{enumerate}

\subsubsection{8.3 Calcul d\textquotesingle un
scoring}\label{83-calcul-dun-scoring}

\begin{longtable}[]{@{}
  >{\raggedright\arraybackslash}p{(\linewidth - 2\tabcolsep) * \real{0.2700}}
  >{\raggedright\arraybackslash}p{(\linewidth - 2\tabcolsep) * \real{0.7000}}@{}}
\toprule\noalign{}
\begin{minipage}[b]{\linewidth}\raggedright
Étape
\end{minipage} & \begin{minipage}[b]{\linewidth}\raggedright
Formule
\end{minipage} \\
\midrule\noalign{}
\endhead
\bottomrule\noalign{}
\endlastfoot
\textbf{Probabilité de défaut} & Sortie de la régression logistique sur
les variables standardisées \\
\textbf{Score} &
\texttt{score\ =\ 515,06\ −\ 28,85\ ×\ ln(\ PD\ ÷\ (1\ −\ PD)\ )}, borné
entre 300 et 850, arrondi à l\textquotesingle entier (les constantes de
CREDIX) \\
\textbf{Décision} & PD \textless{} seuil d\textquotesingle accord :
\textbf{ACCORDÉ} ; PD \textgreater{} seuil de refus : \textbf{REFUSÉ} ;
sinon \textbf{REVUE} \\
\textbf{Contribution d\textquotesingle une variable} &
\texttt{coefficient\ ×\ valeur\ standardisée} \\
\textbf{Trois facteurs} & Les trois variables dont la contribution est
la plus grande \textbf{en valeur absolue} ; le sens est
AUGMENTE\_LE\_RISQUE si la contribution est positive,
DIMINUE\_LE\_RISQUE sinon \\
\end{longtable}

\subsubsection{8.4 Valeurs de
référence}\label{84-valeurs-de-ruxe9fuxe9rence}

Avec le jeu de données de l\textquotesingle annexe A (graine 42), un
premier essai a donné une \textbf{AUC de test d\textquotesingle environ
0,75}, un taux de défaut d\textquotesingle environ 8 \% et, sur les 3
000 clients, environ 78 \% d\textquotesingle accords, 17 \% de revues et
5 \% de refus. Les valeurs exactes dépendent de
l\textquotesingle implémentation ; \textbf{les critères de recette de la
section 12 sont donc écrits avec des marges.}

Trois profils d\textquotesingle exemple, avec le résultat
\textbf{indicatif} obtenu :

\begin{longtable}[]{@{}llllllllll@{}}
\toprule\noalign{}
Profil & Âge & Revenu & Ancienneté & Montant & Durée & Incidents & PD &
Score & Décision \\
\midrule\noalign{}
\endhead
\bottomrule\noalign{}
\endlastfoot
\textbf{A} solide & 42 & 450 000 & 120 mois & 1 500 000 & 24 mois & 0 &
0,5 \% & 667 & ACCORDÉ \\
\textbf{B} intermédiaire & 33 & 220 000 & 36 mois & 2 600 000 & 24 mois
& 2 & 24,5 \% & 548 & REVUE \\
\textbf{C} fragile & 24 & 120 000 & 6 mois & 1 800 000 & 12 mois & 3 &
95 \% & 429 & REFUSÉ \\
\end{longtable}

\begin{center}\rule{0.5\linewidth}{0.5pt}\end{center}

\subsection{9. Interface de programmation
(API)}\label{9-interface-de-programmation-api}

Format des échanges : \textbf{JSON}, avec des noms en
\texttt{snake\_case}. Toutes les routes, sauf \texttt{/health} et
\texttt{/\allowbreak{}api/\allowbreak{}auth/\allowbreak{}login}, exigent
l\textquotesingle en-tête
\texttt{Authorization:\ Bearer\ \textless{}jeton\textgreater{}}.

\begin{longtable}[]{@{}
  >{\raggedright\arraybackslash}p{(\linewidth - 6\tabcolsep) * \real{0.3500}}
  >{\raggedright\arraybackslash}p{(\linewidth - 6\tabcolsep) * \real{0.1220}}
  >{\raggedright\arraybackslash}p{(\linewidth - 6\tabcolsep) * \real{0.3151}}
  >{\raggedright\arraybackslash}p{(\linewidth - 6\tabcolsep) * \real{0.1830}}@{}}
\toprule\noalign{}
\begin{minipage}[b]{\linewidth}\raggedright
Méthode et route
\end{minipage} & \begin{minipage}[b]{\linewidth}\raggedright
Rôle requis
\end{minipage} & \begin{minipage}[b]{\linewidth}\raggedright
Description
\end{minipage} & \begin{minipage}[b]{\linewidth}\raggedright
Codes
\end{minipage} \\
\midrule\noalign{}
\endhead
\bottomrule\noalign{}
\endlastfoot
\texttt{GET\ /health} & aucun & État du système et du modèle & 200 \\
\texttt{POST\ /api/auth/login} & aucun & Connexion & 200, 401 \\
\texttt{POST\ /api/clients} & Agent, Admin & Créer un client & 201,
422 \\
\texttt{GET\ /api/clients?recherche=} & Agent, Admin & Rechercher (nom
ou prénom) & 200 \\
\texttt{GET\ /api/clients/\{id\}} & Agent, Admin & Fiche
d\textquotesingle un client & 200, 404 \\
\texttt{POST\ /api/scorings} & Agent, Admin & Demander un scoring & 201,
404, 422, 503 \\
\texttt{GET\ /api/scorings?client\_id=} & Agent, Admin & Historique
d\textquotesingle un client & 200 \\
\texttt{POST\ /api/admin/utilisateurs} & Admin & Créer un compte & 201,
403, 409, 422 \\
\texttt{GET\ /api/admin/seuils} & Admin & Seuils actifs & 200, 403 \\
\texttt{PUT\ /api/admin/seuils} & Admin & Nouvelle version des seuils &
200, 403, 422 \\
\end{longtable}

Sans jeton ou avec un jeton invalide : \textbf{401}. Avec un rôle
insuffisant : \textbf{403}.

\textbf{Exemple : connexion}

\begin{Verbatim}[breaklines=true,breakanywhere=true,fontsize=\footnotesize,frame=leftline,framerule=1.6pt,rulecolor=\color{credixblue},framesep=3mm,xleftmargin=2mm,xrightmargin=1mm]
POST /api/auth/login
{ "identifiant": "agent1", "mot_de_passe": "MotDePasse2026" }

200 OK
{ "jeton": "eyJhbGciOi...", "role": "AGENT", "expire_a": "2026-09-20T09:00:00Z" }
\end{Verbatim}

\textbf{Exemple : demande de scoring}

\begin{Verbatim}[breaklines=true,breakanywhere=true,fontsize=\footnotesize,frame=leftline,framerule=1.6pt,rulecolor=\color{credixblue},framesep=3mm,xleftmargin=2mm,xrightmargin=1mm]
POST /api/scorings
{ "client_id": 12, "montant_credit": 2600000, "duree_mois": 24 }

201 Created
{
  "id": 57,
  "client_id": 12,
  "probabilite_defaut": 0.245,
  "score": 548,
  "decision": "REVUE",
  "facteurs": [
    { "libelle": "Incidents de paiement passés", "contribution": 0.931, "sens": "AUGMENTE_LE_RISQUE" },
    { "libelle": "Taux d'effort", "contribution": 0.408, "sens": "AUGMENTE_LE_RISQUE" },
    { "libelle": "Âge", "contribution": 0.256, "sens": "AUGMENTE_LE_RISQUE" }
  ],
  "modele_version": "1.0",
  "seuil_accord": 0.10,
  "seuil_refus": 0.30,
  "date_heure": "2026-09-19T14:03:11Z"
}
\end{Verbatim}

\begin{center}\rule{0.5\linewidth}{0.5pt}\end{center}

\subsection{10. Interface utilisateur}\label{10-interface-utilisateur}

Une \textbf{page web unique}, sans compilation, servie à la racine
\texttt{/}. Trois zones, affichées selon le rôle.

\begin{longtable}[]{@{}
  >{\raggedright\arraybackslash}p{(\linewidth - 2\tabcolsep) * \real{0.3233}}
  >{\raggedright\arraybackslash}p{(\linewidth - 2\tabcolsep) * \real{0.6467}}@{}}
\toprule\noalign{}
\begin{minipage}[b]{\linewidth}\raggedright
Zone
\end{minipage} & \begin{minipage}[b]{\linewidth}\raggedright
Contenu
\end{minipage} \\
\midrule\noalign{}
\endhead
\bottomrule\noalign{}
\endlastfoot
\textbf{Connexion} & Identifiant, mot de passe, bouton « Se connecter »
; message clair en cas d\textquotesingle échec \\
\textbf{Clients et scoring} (Agent et Admin) & Barre de recherche ;
formulaire « Nouveau client » ; fiche d\textquotesingle un client avec
le formulaire de scoring (montant, durée) ; \textbf{résultat} : score,
décision (colorée : vert, orange, rouge), probabilité de défaut, trois
facteurs avec leur sens ; \textbf{historique} du client \\
\textbf{Administration} (Admin seulement) & Création
d\textquotesingle un compte ; réglage des deux seuils, avec le rappel de
la règle « accord \textless{} refus » \\
\end{longtable}

Un bouton « Se déconnecter » efface le jeton.

\begin{center}\rule{0.5\linewidth}{0.5pt}\end{center}

\subsection{11. Plan de réalisation par
modules}\label{11-plan-de-ruxe9alisation-par-modules}

L\textquotesingle ordre va du plus indépendant au plus dépendant. Chaque
module se termine quand ses tests passent.

\begin{longtable}[]{@{}
  >{\raggedright\arraybackslash}p{(\linewidth - 6\tabcolsep) * \real{0.1750}}
  >{\raggedright\arraybackslash}p{(\linewidth - 6\tabcolsep) * \real{0.2850}}
  >{\raggedright\arraybackslash}p{(\linewidth - 6\tabcolsep) * \real{0.2250}}
  >{\raggedright\arraybackslash}p{(\linewidth - 6\tabcolsep) * \real{0.2850}}@{}}
\toprule\noalign{}
\begin{minipage}[b]{\linewidth}\raggedright
Module
\end{minipage} & \begin{minipage}[b]{\linewidth}\raggedright
Objectif
\end{minipage} & \begin{minipage}[b]{\linewidth}\raggedright
Fichiers principaux
\end{minipage} & \begin{minipage}[b]{\linewidth}\raggedright
Terminé quand
\end{minipage} \\
\midrule\noalign{}
\endhead
\bottomrule\noalign{}
\endlastfoot
\textbf{M1 Socle} & Structure du projet, configuration, base et tables &
\texttt{app/config.py}, \texttt{database.py}, \texttt{models.py},
\texttt{requirements.\allowbreak{}txt} & Les tables sont créées ; test
de connexion à la base \\
\textbf{M2 Données et modèle} & Générer les données, entraîner et
sauvegarder le modèle (\textbf{EF-10}, ENF-3) &
\texttt{scripts/\allowbreak{}generer\_\allowbreak{}donnees.\allowbreak{}py},
\texttt{scripts/\allowbreak{}entrainer\_\allowbreak{}modele.\allowbreak{}py}
& AUC de test au moins égale à 0,70 ; deux exécutions donnent les mêmes
coefficients \\
\textbf{M3 Service de scoring} & PD, score, décision, trois facteurs
(\textbf{EF-5, EF-6}) &
\texttt{app/\allowbreak{}services/\allowbreak{}scoring.\allowbreak{}py}
& Les profils A et C donnent ACCORDÉ et REFUSÉ ; plus
d\textquotesingle incidents implique une PD plus élevée \\
\textbf{M4 Sécurité et comptes} & Hachage, jeton, rôles (\textbf{EF-1,
EF-2, EF-8}) & \texttt{app/\allowbreak{}security.\allowbreak{}py},
\texttt{routers/\allowbreak{}auth.\allowbreak{}py},
\texttt{routers/\allowbreak{}admin.\allowbreak{}py},
\texttt{scripts/\allowbreak{}creer\_\allowbreak{}admin.\allowbreak{}py}
& Connexion réussie ou refusée ; 403 pour un Agent sur une route
d\textquotesingle administration \\
\textbf{M5 API métier} & Clients, scorings, historique, seuils
(\textbf{EF-3, EF-4, EF-7, EF-9}, ENF-1) &
\texttt{routers/\allowbreak{}clients.\allowbreak{}py},
\texttt{routers/\allowbreak{}scorings.\allowbreak{}py},
\texttt{services/\allowbreak{}seuils.\allowbreak{}py} & Tous les tests
d\textquotesingle API passent \\
\textbf{M6 Interface et livraison} & Page web, README de lancement
(ENF-4) & \texttt{frontend/*}, \texttt{README.md} & Le README est suivi
à la lettre sur une machine neuve \\
\end{longtable}

\textbf{Dépendances :} M1 avant tout ; M2 avant M3 ; M4 avant M5 ; M3 et
M5 avant M6.

\textbf{Règles de conduite pour Claude Code :} un seul module à la fois
; les tests d\textquotesingle abord ; ne jamais modifier un module déjà
validé sans le dire ; ne jamais écrire de mot de passe ni de secret dans
le code.

\begin{center}\rule{0.5\linewidth}{0.5pt}\end{center}

\subsection{12. Plan de tests et
recette}\label{12-plan-de-tests-et-recette}

\subsubsection{12.1 Tests automatiques}\label{121-tests-automatiques}

Une commande, \texttt{pytest}, lance tous les tests. Il y a \textbf{au
moins un test par exigence fonctionnelle}, plus :

\begin{itemize}
\tightlist
\item
  \textbf{monotonie} : à profil identique, la PD \textbf{augmente} avec
  le nombre d\textquotesingle incidents (0, 1, 2, 3) ;
\item
  \textbf{bornes} : le score est toujours entre 300 et 850 ;
\item
  \textbf{traçabilité} : un scoring enregistré contient la version du
  modèle et les seuils appliqués ;
\item
  \textbf{sécurité} : le hachage du mot de passe est différent du mot de
  passe, et la vérification réussit.
\end{itemize}

\subsubsection{12.2 Scénario de recette (à la
main)}\label{122-scuxe9nario-de-recette-uxe0-la-main}

\begin{longtable}[]{@{}
  >{\raggedright\arraybackslash}p{(\linewidth - 4\tabcolsep) * \real{0.0875}}
  >{\raggedright\arraybackslash}p{(\linewidth - 4\tabcolsep) * \real{0.4412}}
  >{\raggedright\arraybackslash}p{(\linewidth - 4\tabcolsep) * \real{0.4412}}@{}}
\toprule\noalign{}
\begin{minipage}[b]{\linewidth}\raggedright
N°
\end{minipage} & \begin{minipage}[b]{\linewidth}\raggedright
Action
\end{minipage} & \begin{minipage}[b]{\linewidth}\raggedright
Résultat attendu
\end{minipage} \\
\midrule\noalign{}
\endhead
\bottomrule\noalign{}
\endlastfoot
1 & Lancer les commandes du README & L\textquotesingle application
répond sur \texttt{http:/\allowbreak{}/\allowbreak{}localhost:8000} ;
\texttt{/health} indique « ok » \\
2 & Se connecter avec le compte administrateur & Accès à
l\textquotesingle écran d\textquotesingle administration \\
3 & Créer un compte « agent1 » & Le compte apparaît ; il peut se
connecter \\
4 & Se connecter en agent1 & Pas d\textquotesingle accès à
l\textquotesingle administration \\
5 & Créer le client du profil \textbf{A} & Le client est retrouvable par
recherche \\
6 & Demander un scoring (1 500 000 sur 24 mois) & Décision
\textbf{ACCORDÉ}, score élevé, trois facteurs \\
7 & Créer le client du profil \textbf{C} et demander un scoring (1 800
000 sur 12 mois) & Décision \textbf{REFUSÉ} \\
8 & Consulter l\textquotesingle historique du client \textbf{A} & Le
scoring figure, avec version du modèle et seuils \\
9 & En administrateur, passer le seuil de refus de 0,30 à 0,20, puis
rejouer le scoring du profil \textbf{B} (probabilité de défaut
d\textquotesingle environ 24,5 \%) & La décision passe de \textbf{REVUE
à REFUSÉ} ; l\textquotesingle ancien scoring \textbf{garde} ses seuils
d\textquotesingle origine (0,30) \\
10 & Essayer un seuil d\textquotesingle accord de 0,40 avec un seuil de
refus de 0,20 & Le système refuse (422) et explique que
l\textquotesingle accord doit être inférieur au refus \\
\end{longtable}

\subsubsection{12.3 Critères de succès de la
démonstration}\label{123-crituxe8res-de-succuxe8s-de-la-duxe9monstration}

\begin{longtable}[]{@{}
  >{\raggedright\arraybackslash}p{(\linewidth - 2\tabcolsep) * \real{0.6117}}
  >{\raggedright\arraybackslash}p{(\linewidth - 2\tabcolsep) * \real{0.3583}}@{}}
\toprule\noalign{}
\begin{minipage}[b]{\linewidth}\raggedright
Résultat
\end{minipage} & \begin{minipage}[b]{\linewidth}\raggedright
Signification
\end{minipage} \\
\midrule\noalign{}
\endhead
\bottomrule\noalign{}
\endlastfoot
Le \textbf{plan} produit par Claude Code couvre tous les modules M1 à
M6, dans un ordre cohérent & Le document de conception est
\textbf{exploitable} \\
Le module \textbf{M2} livre une AUC d\textquotesingle au moins 0,70 & Le
premier module tient ses promesses \\
Les tests du module \textbf{M3} passent & La logique métier est
correcte \\
\end{longtable}

\begin{center}\rule{0.5\linewidth}{0.5pt}\end{center}

\subsection{13. Limites et évolution vers
CREDIX}\label{13-limites-et-uxe9volution-vers-credix}

\begin{longtable}[]{@{}
  >{\raggedright\arraybackslash}p{(\linewidth - 4\tabcolsep) * \real{0.3706}}
  >{\raggedright\arraybackslash}p{(\linewidth - 4\tabcolsep) * \real{0.2180}}
  >{\raggedright\arraybackslash}p{(\linewidth - 4\tabcolsep) * \real{0.3815}}@{}}
\toprule\noalign{}
\begin{minipage}[b]{\linewidth}\raggedright
Ce que fait CREDIX
\end{minipage} & \begin{minipage}[b]{\linewidth}\raggedright
Le MVP
\end{minipage} & \begin{minipage}[b]{\linewidth}\raggedright
Comment y arriver
\end{minipage} \\
\midrule\noalign{}
\endhead
\bottomrule\noalign{}
\endlastfoot
Deux flux : score et détection des profils atypiques (autoencodeur) & Un
seul flux & Ajouter un second modèle non supervisé qui peut
\textbf{durcir} une décision \\
WoE, sélection de variables, LightGBM & Régression logistique sur six
variables & Remplacer le modèle derrière la même interface de scoring \\
Calibration isotonique des probabilités & Aucune & Ajouter une étape
entre le modèle et le score \\
Indice de couverture ρc et clients à dossier incomplet & Aucun &
Calculer la part d\textquotesingle information disponible par client \\
Explications par SHAP, phrases en langage naturel & Contributions du
modèle linéaire & Remplacer par SHAP et un gabarit de phrases \\
Rôle \textbf{Superviseur} et file de revue & La décision REVUE
n\textquotesingle a pas de suite & Ajouter le rôle, la file, la décision
justifiée \\
MongoDB, Docker, service PDF, journal d\textquotesingle audit & SQLite,
aucun conteneur & Voir l\textquotesingle architecture de CREDIX \\
\end{longtable}

Le MVP n\textquotesingle a \textbf{aucune valeur prédictive réelle} :
ses données sont synthétiques. Il sert à \textbf{valider la méthode},
pas à évaluer un risque.

\begin{center}\rule{0.5\linewidth}{0.5pt}\end{center}

\subsection{14. Glossaire}\label{14-glossaire}

\begin{longtable}[]{@{}
  >{\raggedright\arraybackslash}p{(\linewidth - 2\tabcolsep) * \real{0.2771}}
  >{\raggedright\arraybackslash}p{(\linewidth - 2\tabcolsep) * \real{0.6929}}@{}}
\toprule\noalign{}
\begin{minipage}[b]{\linewidth}\raggedright
Terme
\end{minipage} & \begin{minipage}[b]{\linewidth}\raggedright
Définition
\end{minipage} \\
\midrule\noalign{}
\endhead
\bottomrule\noalign{}
\endlastfoot
\textbf{Probabilité de défaut (PD)} & Chance estimée
qu\textquotesingle un client ne rembourse pas son crédit \\
\textbf{Score} & La PD convertie en un nombre de 300 à 850 : plus il est
haut, plus le risque est faible \\
\textbf{Décision} & ACCORDÉ, REFUSÉ ou REVUE (à trancher par un
humain) \\
\textbf{Seuils} & Les deux probabilités qui séparent les trois
décisions \\
\textbf{Taux d\textquotesingle effort} & Part du revenu mensuel
consacrée au remboursement du crédit \\
\textbf{AUC} & Mesure de la capacité d\textquotesingle un modèle à
séparer bons et mauvais payeurs (0,5 : aucun pouvoir ; 1 : parfait) \\
\textbf{Régression logistique} & Modèle statistique simple qui calcule
une probabilité à partir d\textquotesingle une somme pondérée de
variables \\
\textbf{Variable standardisée} & Variable ramenée à une moyenne de 0 et
un écart-type de 1 \\
\textbf{JWT (jeton)} & Justificatif temporaire délivré après
connexion \\
\textbf{PBKDF2} & Méthode de hachage de mot de passe \\
\textbf{Recette} & Vérification finale, à la main, que le système répond
au besoin \\
\end{longtable}

\begin{center}\rule{0.5\linewidth}{0.5pt}\end{center}

\subsection{15. Annexes}\label{15-annexes}

\subsubsection{Annexe A : générateur de données synthétiques
(référence)}\label{annexe-a--guxe9nuxe9rateur-de-donnuxe9es-synthuxe9tiques-ruxe9fuxe9rence}

Ce script définit le « monde » du MVP. Le donner tel quel à Claude Code
garantit des données \textbf{identiques d\textquotesingle une exécution
à l\textquotesingle autre} (graine 42).

\begin{Verbatim}[breaklines=true,breakanywhere=true,fontsize=\footnotesize,frame=leftline,framerule=1.6pt,rulecolor=\color{credixblue},framesep=3mm,xleftmargin=2mm,xrightmargin=1mm]
import numpy as np
import pandas as pd

rng = np.random.default_rng(42)
n = 3000
age = np.clip(rng.normal(38, 10, n), 21, 65).round().astype(int)
revenu = np.clip(np.exp(rng.normal(np.log(250000), 0.5, n)), 80000, 2000000).round(-3)
anciennete = np.clip(rng.exponential(60, n), 0, 360).round().astype(int)
duree = rng.choice([6, 12, 18, 24, 36, 48], n, p=[.10, .25, .20, .25, .15, .05])
montant = (revenu * rng.uniform(1, 10, n)).round(-4)
incidents = np.clip(rng.poisson(0.6, n), 0, 8)
taux_effort = (montant / duree) / revenu

logit = (-4.4 - 0.035 * (age - 38) - 1.1 * (np.log(revenu) - np.log(250000))
         - 0.004 * (anciennete - 60) + 3.0 * taux_effort
         + 0.02 * (duree - 18) + 0.55 * incidents)
probabilite = 1 / (1 + np.exp(-logit))
defaut = rng.binomial(1, probabilite)

donnees = pd.DataFrame({
    "age": age, "revenu_mensuel": revenu, "anciennete_emploi_mois": anciennete,
    "montant_credit": montant, "duree_mois": duree,
    "nb_incidents_paiement": incidents, "defaut": defaut,
})
donnees.to_csv("clients_synthetiques.csv", index=False)
\end{Verbatim}

\subsubsection{Annexe B : prompt de
démonstration}\label{annexe-b--prompt-de-duxe9monstration}

À donner à Claude Code, en \textbf{mode Plan}, dans un dossier contenant
ce document sous
\texttt{docs/\allowbreak{}conception\_\allowbreak{}mvp.\allowbreak{}md}.

\begin{Verbatim}[breaklines=true,breakanywhere=true,fontsize=\footnotesize,frame=leftline,framerule=1.6pt,rulecolor=\color{credixblue},framesep=3mm,xleftmargin=2mm,xrightmargin=1mm]
/plan Lis docs/conception_mvp.md en entier. Propose un plan de réalisation par
modules (M1 à M6 ou mieux, si tu vois une meilleure découpe). Pour chaque
module : l'objectif, les fichiers à créer, les tests à écrire, le critère qui
dit qu'il est terminé, les dépendances. Signale les points du document qui te
paraissent ambigus ou contradictoires. Ne crée aucun fichier pour l'instant.
\end{Verbatim}

Puis, après relecture et approbation du plan :

\begin{Verbatim}[breaklines=true,breakanywhere=true,fontsize=\footnotesize,frame=leftline,framerule=1.6pt,rulecolor=\color{credixblue},framesep=3mm,xleftmargin=2mm,xrightmargin=1mm]
Réalise le module M2 du plan. Écris d'abord les tests, puis le code. Lance les
tests et montre-moi le résultat, dont l'AUC de test. Ne touche à aucun autre module.
\end{Verbatim}

\subsubsection{Annexe C : commandes de lancement attendues à la
fin}\label{annexe-c--commandes-de-lancement-attendues-uxe0-la-fin}

\begin{Verbatim}[breaklines=true,breakanywhere=true,fontsize=\footnotesize,frame=leftline,framerule=1.6pt,rulecolor=\color{credixblue},framesep=3mm,xleftmargin=2mm,xrightmargin=1mm]
python3 -m venv .venv && source .venv/bin/activate
pip install -r requirements.txt
cp .env.example .env
python scripts/generer_donnees.py && python scripts/entrainer_modele.py
python scripts/creer_admin.py
uvicorn app.main:app --port 8000
\end{Verbatim}

\subsubsection{Annexe D : ce que ce document ne dit pas
volontairement}\label{annexe-d--ce-que-ce-document-ne-dit-pas-volontairement}

Le \textbf{code} des modules M3 à M6, la structure exacte des tests et
l\textquotesingle aspect visuel de la page. C\textquotesingle est le
travail de Claude Code, et c\textquotesingle est ce que la démonstration
cherche à observer.

\end{document}
